\begin{quote}\small Article 219 (Modifié par la loi n° 64-34 du 2 juillet 1964 et par la loi n° 1989-23 du 27 février 1989).- Quand les violences ci-dessus exprimées ont été suivies de mutilation, perte de l'usage d'un membre, défiguration, infirmité ou incapacité permanente dont le taux ne dépasse pas 20\%, le coupable est puni de cinq ans d'emprisonnement. La peine sera de dix ans de prison, s'il est résulté de ces sortes de violence une incapacité dont le taux dépasse 20\%. La peine est portée à douze ans d'emprisonnement si le coupable est un descendant de la victime, quel que soit le taux de l'incapacité, même en cas de désistement.\end{quote}